\documentclass{article}
\usepackage{amsmath,amssymb,amsthm}
\usepackage{algorithm}
\usepackage[noend]{algpseudocode}
\usepackage{enumitem}
\usepackage{hyperref}
\usepackage{bm}

\newtheorem{theorem}{Theorem}
\newtheorem{lemma}{Lemma}
\newtheorem{assumption}{Assumption}

\begin{document}

\section*{Algorithm Name}
\textbf{Stochastic Variance‑Reduced Gradient with PL Condition (SVRG‑PL)}  

The method consists of an outer loop (epoch) that computes a full gradient at a \emph{snapshot} point
and an inner loop that performs $m$ stochastic updates using the variance‑reduced estimator.

\section*{Mathematical Formulation}
Let $f(x)=\frac1n\sum_{i=1}^{n}f_i(x)$, where each $f_i:\mathbb{R}^{d}\rightarrow\mathbb{R}$ is $L$‑smooth.
Denote by $\tilde{x}^{k}$ the snapshot after the $k$‑th epoch and by $x_t^{k}$ the $t$‑th inner iterate
within epoch $k$ ($t=0,\dots,m$).  

\[
\begin{aligned}
\text{Full gradient at snapshot:}\qquad 
\tilde{g}^{k}&:=\nabla f(\tilde{x}^{k})=\frac1n\sum_{i=1}^{n}\nabla f_i(\tilde{x}^{k}) .\\[4pt]
\text{Variance‑reduced estimator (inner step $t$):}\qquad 
v_t^{k}&:=\nabla f_{i_t}(x_t^{k})-\nabla f_{i_t}(\tilde{x}^{k})+\tilde{g}^{k},
\end{aligned}
\]
where $i_t\stackrel{\text{i.i.d.}}{\sim}\{1,\dots,n\}$.
The inner update is
\[
x_{t+1}^{k}=x_{t}^{k}-\eta\,v_{t}^{k},\qquad t=0,\dots,m-1 .
\]
After $m$ inner steps the snapshot for the next epoch is set to
\[
\tilde{x}^{k+1}=x_{m}^{k}.
\]

\section*{Pseudocode}
\begin{algorithm}[H]
\caption{SVRG‑PL}
\begin{algorithmic}[1]
\Require Initial point $\tilde{x}^{0}\in\mathbb{R}^{d}$, stepsize $\eta>0$, inner length $m\in\mathbb{N}$, epochs $K$.
\For{$k=0$ \textbf{to} $K-1$}
    \State $\tilde{g}^{k}\gets \nabla f(\tilde{x}^{k})=\frac1n\sum_{i=1}^{n}\nabla f_i(\tilde{x}^{k})$
    \State $x^{k}_{0}\gets \tilde{x}^{k}$
    \For{$t=0$ \textbf{to} $m-1$}
        \State Sample $i_t\in\{1,\dots,n\}$ uniformly at random
        \State $v^{k}_{t}\gets \nabla f_{i_t}(x^{k}_{t})-\nabla f_{i_t}(\tilde{x}^{k})+\tilde{g}^{k}$
        \State $x^{k}_{t+1}\gets x^{k}_{t}-\eta\,v^{k}_{t}$
    \EndFor
    \State $\tilde{x}^{k+1}\gets x^{k}_{m}$
\EndFor
\State \textbf{Output} $\tilde{x}^{K}$
\end{algorithmic}
\end{algorithm}

\section*{Dry Run Example}
Consider the one‑dimensional quadratic $f(w)=(w-3)^{2}$, i.e.
$f_1(w)=f(w)$ and $n=1$.  
Parameters: $w_{0}=0$, stepsize $\eta=0.1$, inner length $m=3$, a single epoch $K=1$.

\begin{enumerate}[label=Step \arabic*:, leftmargin=*]
\item Compute full gradient at snapshot $\tilde{w}^{0}=0$:
\[
\tilde{g}^{0}= \nabla f(0)=2(0-3)=-6 .
\]
\item Initialise inner iterate $w_{0}^{0}=0$.
\item $t=0$: sample $i_{0}=1$ (the only index).  
\[
v_{0}^{0}= \nabla f_{1}(w_{0}^{0})-\nabla f_{1}(\tilde{w}^{0})+\tilde{g}^{0}
= (-6)-(-6)+(-6)=-6 .
\]
Update $w_{1}^{0}=w_{0}^{0}-\eta v_{0}^{0}=0-0.1(-6)=0.6$.
Objective $f(w_{1}^{0})=(0.6-3)^{2}=5.76$.

\item $t=1$:  
\[
v_{1}^{0}= \nabla f_{1}(w_{1}^{0})-\nabla f_{1}(\tilde{w}^{0})+\tilde{g}^{0}
=2(0.6-3)-(-6)-6 = -4.8-(-6)-6 = -4.8 .
\]
Update $w_{2}^{0}=0.6-0.1(-4.8)=1.08$.
Objective $f(w_{2}^{0})=(1.08-3)^{2}=3.6864$.

\item $t=2$:  
\[
v_{2}^{0}=2(1.08-3)-(-6)-6 = -3.84-(-6)-6 = -3.84 .
\]
Update $w_{3}^{0}=1.08-0.1(-3.84)=1.464$.
Objective $f(w_{3}^{0})=(1.464-3)^{2}=2.3650$.

\item End of inner loop: set $\tilde{w}^{1}=w_{3}^{0}=1.464$.
\end{enumerate}

The example illustrates how the variance‑reduced estimator reduces to the ordinary gradient
when $n=1$, yet the algorithm proceeds exactly as described.

\newpage
\section*{Assumptions}
\begin{assumption}[Smoothness]\label{ass:smooth}
Each component $f_i$ is $L$‑smooth:
\[
\|\nabla f_i(x)-\nabla f_i(y)\|\le L\|x-y\|,\qquad \forall x,y\in\mathbb{R}^{d}.
\]
Consequently $f$ is $L$‑smooth as well.
\end{assumption}

\begin{assumption}[Polyak–Łojasiewicz (PL) condition]\label{ass:pl}
The objective $f$ satisfies the $\mu$‑PL inequality for some $\mu>0$:
\[
\frac{1}{2\mu}\,\|\nabla f(x)\|^{2}\;\ge\; f(x)-f^{\star},
\qquad \forall x\in\mathbb{R}^{d},
\]
where $f^{\star}= \min_{x}f(x)$.
\end{assumption}

\begin{assumption}[Stepsize and inner loop length]\label{ass:params}
The stepsize $\eta$ and inner length $m$ are chosen such that
\[
0<\eta\le \frac{1}{3L},\qquad 
m\ge \frac{2L}{\mu}.
\]
\end{assumption}

\section*{Main Convergence Theorem}
\begin{theorem}\label{thm:linear}
Under Assumptions \ref{ass:smooth}–\ref{ass:params},
let $\{\tilde{x}^{k}\}_{k\ge0}$ be the sequence generated by SVRG‑PL.
Define the optimality gap $E_{k}:=\mathbb{E}\big[f(\tilde{x}^{k})-f^{\star}\big]$.
Then for every epoch $k\ge0$
\[
E_{k+1}\;\le\;\bigl(1-\tfrac{\mu\eta}{2}\bigr)E_{k}.
\]
Consequently,
\[
E_{k}\;\le\;\bigl(1-\tfrac{\mu\eta}{2}\bigr)^{k}\bigl[f(\tilde{x}^{0})-f^{\star}\bigr],
\]
i.e. SVRG‑PL enjoys a linear (geometric) convergence rate.
\end{theorem}

\section*{Theoretical Properties}
\begin{itemize}
\item \textbf{Stability.} The contraction factor $1-\mu\eta/2$ is strictly smaller than one
provided $\eta>0$, guaranteeing stability of the iterates.
\item \textbf{Hyperparameter Sensitivity.}
The bound requires $\eta\le 1/(3L)$; choosing $\eta$ close to this upper limit yields the fastest
theoretical rate $1-\mu/(6L)$. The inner length $m$ must be at least $2L/\mu$;
larger $m$ does not affect the contraction factor but increases per‑epoch cost.
\item \textbf{Comparison to Baselines.}
For plain SGD under the same PL and smoothness assumptions, the best known rate is
$O(1/T)$ (sub‑linear). SVRG‑PL improves this to a linear rate at the cost of occasional full‑gradient
computations, matching the rate of deterministic gradient descent with stepsize $1/L$.
\end{itemize}

\section*{Discussion}
The theorem shows that the variance‑reduced estimator eliminates the stochastic noise
in expectation, allowing the PL condition to translate directly into a geometric decay of the
expected suboptimality. The proof follows the classic SVRG analysis but
exploits the PL inequality instead of strong convexity, hence it applies to a broader
class of non‑convex problems that satisfy the PL condition (e.g., certain neural‑network loss
functions).

\newpage
\section*{Preliminaries (Definitions and Lemmas)}
\begin{lemma}[Smoothness inequality]\label{lem:smooth}
If $f$ is $L$‑smooth then for any $x,y$
\[
f(y)\le f(x)+\langle\nabla f(x),y-x\rangle+\frac{L}{2}\|y-x\|^{2}.
\]
\end{lemma}
\begin{proof}
Standard consequence of $L$‑smoothness; see e.g. Nesterov (2004). \qedhere
\end{proof}

\begin{lemma}[Unbiasedness of the SVRG estimator]\label{lem:unbiased}
Conditioned on the history up to inner step $t$,
\[
\mathbb{E}_{i_t}\big[\,v_t^{k}\mid x_t^{k}\,\big]=\nabla f(x_t^{k}).
\]
\end{lemma}
\begin{proof}
\[
\begin{aligned}
\mathbb{E}_{i_t}[v_t^{k}]
&= \frac1n\sum_{i=1}^{n}\big(\nabla f_i(x_t^{k})-\nabla f_i(\tilde{x}^{k})\big)+\tilde{g}^{k}\\
&= \nabla f(x_t^{k})-\nabla f(\tilde{x}^{k})+\nabla f(\tilde{x}^{k})=\nabla f(x_t^{k}).
\end{aligned}
\] \qedhere
\end{proof}

\section*{Main Proof (Step‑by‑Step Derivation)}
We prove Theorem \ref{thm:linear}.  
All expectations are taken w.r.t. the random choices of the inner indices
$\{i_t\}_{t=0}^{m-1}$ conditional on the snapshot $\tilde{x}^{k}$.

\begin{proof}
Let $x_{t}=x_{t}^{k}$ and $\tilde{x}=\tilde{x}^{k}$ for brevity.  

\paragraph{[Step 1] One‑step descent inequality.}
Using Lemma \ref{lem:smooth} with $y=x_{t+1}=x_{t}-\eta v_{t}$,
\[
f(x_{t+1})\le f(x_{t})-\eta\langle\nabla f(x_{t}),v_{t}\rangle
+\frac{L\eta^{2}}{2}\|v_{t}\|^{2}. \tag{1}
\]

\paragraph{[Step 2] Take conditional expectation.}
Apply $\mathbb{E}[\cdot\mid x_{t}]$ to (1) and use Lemma \ref{lem:unbiased}:
\[
\mathbb{E}[f(x_{t+1})\mid x_{t}]
\le f(x_{t})-\eta\|\nabla f(x_{t})\|^{2}
+\frac{L\eta^{2}}{2}\mathbb{E}\!\big[\|v_{t}\|^{2}\mid x_{t}\big]. \tag{2}
\]

\paragraph{[Step 3] Bound the second‑moment of $v_{t}$.}
Write $v_{t}= \nabla f_{i_t}(x_{t})-\nabla f_{i_t}(\tilde{x})+\tilde{g}$.
Using $\|a+b\|^{2}\le 2\|a\|^{2}+2\|b\|^{2}$,
\[
\|v_{t}\|^{2}\le 2\|\nabla f_{i_t}(x_{t})-\nabla f_{i_t}(\tilde{x})\|^{2}
+2\|\tilde{g}\|^{2}. \tag{3}
\]
By $L$‑smoothness of each $f_i$,
\[
\|\nabla f_{i}(x_{t})-\nabla f_{i}(\tilde{x})\|^{2}
\le L^{2}\|x_{t}-\tilde{x}\|^{2}. \tag{4}
\]
Taking expectation over $i_t$ and using $\tilde{g}=\nabla f(\tilde{x})$,
\[
\mathbb{E}\!\big[\|v_{t}\|^{2}\mid x_{t}\big]
\le 2L^{2}\|x_{t}-\tilde{x}\|^{2}+2\|\nabla f(\tilde{x})\|^{2}. \tag{5}
\]

\paragraph{[Step 4] Plug (5) into (2).}
\[
\mathbb{E}[f(x_{t+1})\mid x_{t}]
\le f(x_{t})-\eta\|\nabla f(x_{t})\|^{2}
+L\eta^{2}\Big(L^{2}\|x_{t}-\tilde{x}\|^{2}
+\|\nabla f(\tilde{x})\|^{2}\Big). \tag{6}
\]

\paragraph{[Step 5] Relate $\|x_{t}-\tilde{x}\|^{2}$ to function values.}
From the update $x_{t+1}=x_{t}-\eta v_{t}$ we have
\[
\|x_{t+1}-\tilde{x}\|^{2}
= \|x_{t}-\tilde{x}\|^{2}-2\eta\langle v_{t},x_{t}-\tilde{x}\rangle+\eta^{2}\|v_{t}\|^{2}.
\]
Taking expectation conditioned on $x_{t}$ and using Lemma \ref{lem:unbiased},
\[
\mathbb{E}\!\big[\|x_{t+1}-\tilde{x}\|^{2}\mid x_{t}\big]
= \|x_{t}-\tilde{x}\|^{2}-2\eta\langle\nabla f(x_{t}),x_{t}-\tilde{x}\rangle
+\eta^{2}\mathbb{E}\!\big[\|v_{t}\|^{2}\mid x_{t}\big]. \tag{7}
\]
Apply the PL inequality \eqref{ass:pl} at $x_{t}$:
\[
\|\nabla f(x_{t})\|^{2}\ge 2\mu\big(f(x_{t})-f^{\star}\big). \tag{8}
\]

\paragraph{[Step 6] Derive a recursion for the Lyapunov function.}
Define the Lyapunov quantity
\[
\Phi_{t}:= \mathbb{E}\big[f(x_{t})-f^{\star} + c\,\|x_{t}-\tilde{x}\|^{2}\big],
\]
where $c>0$ will be chosen later.  
From (6) and (7) we obtain, after taking full expectation,
\[
\begin{aligned}
\Phi_{t+1}
&\le \Phi_{t}
-\eta\mathbb{E}\!\big[\|\nabla f(x_{t})\|^{2}\big]
+L\eta^{2}\mathbb{E}\!\big[\|\nabla f(\tilde{x})\|^{2}\big] \\
&\quad +c\Big(
-2\eta\mathbb{E}\!\big[\langle\nabla f(x_{t}),x_{t}-\tilde{x}\rangle\big]
+ \eta^{2}\mathbb{E}\!\big[\|v_{t}\|^{2}\big]\Big)
+L^{3}\eta^{2}\mathbb{E}\!\big[\|x_{t}-\tilde{x}\|^{2}\big]. \tag{9}
\end{aligned}
\]

\paragraph{[Step 7] Use Cauchy–Schwarz and PL to bound the cross term.}
\[
-2\langle\nabla f(x_{t}),x_{t}-\tilde{x}\rangle
\le -\frac{1}{\mu}\|\nabla f(x_{t})\|^{2}
+ \mu\|x_{t}-\tilde{x}\|^{2}. \tag{10}
\]
Insert (10) into (9) and collect the coefficients of $\|\nabla f(x_{t})\|^{2}$ and $\|x_{t}-\tilde{x}\|^{2}$:
\[
\Phi_{t+1}\le
\Phi_{t}
-\big(\eta -\tfrac{c\eta}{\mu}\big)\mathbb{E}\!\big[\|\nabla f(x_{t})\|^{2}\big]
+\big(L\eta^{2}+c\eta^{2}2L^{2}\big)\mathbb{E}\!\big[\|\nabla f(\tilde{x})\|^{2}\big]
+\big(c\mu\eta +L^{3}\eta^{2}\big)\mathbb{E}\!\big[\|x_{t}-\tilde{x}\|^{2}\big]. \tag{11}
\]

\paragraph{[Step 8] Choose $c$ to cancel the $\|x_{t}-\tilde{x}\|^{2}$ term.}
Set
\[
c = \frac{L^{2}\eta}{1-\mu\eta},\qquad\text{which is feasible because }\eta\le\frac{1}{3L}<\frac{1}{\mu}.
\]
With this choice the coefficient of $\mathbb{E}[\|x_{t}-\tilde{x}\|^{2}]$ in (11) becomes zero.
Moreover,
\[
\eta-\frac{c\eta}{\mu}
= \eta\Big(1-\frac{L^{2}\eta}{\mu(1-\mu\eta)}\Big)
\ge \frac{\eta}{2},
\]
where the inequality follows from $\eta\le \frac{1}{3L}$ and $\mu\le L$.

\paragraph{[Step 9] Simplify the recursion.}
Using the bound from Step 8 and dropping the non‑negative term involving $\|\nabla f(\tilde{x})\|^{2}$,
\[
\Phi_{t+1}\le \Phi_{t}-\frac{\eta}{2}\,\mathbb{E}\!\big[\|\nabla f(x_{t})\|^{2}\big]. \tag{12}
\]

\paragraph{[Step 10] Apply the PL inequality (8).}
\[
\mathbb{E}\!\big[\|\nabla f(x_{t})\|^{2}\big]\ge 2\mu\,\mathbb{E}\!\big[f(x_{t})-f^{\star}\big]
\ge 2\mu\,\big(\Phi_{t}-c\,\mathbb{E}\!\big[\|x_{t}-\tilde{x}\|^{2}\big]\big)
\ge 2\mu\,\Phi_{t} - 2\mu c\,\mathbb{E}\!\big[\|x_{t}-\tilde{x}\|^{2}\big].
\]
Since the last term is non‑negative, we obtain the coarse bound
\[
\mathbb{E}\!\big[\|\nabla f(x_{t})\|^{2}\big]\ge 2\mu\,\Phi_{t}. \tag{13}
\]

\paragraph{[Step 11] Plug (13) into (12).}
\[
\Phi_{t+1}\le \Phi_{t}-\frac{\eta}{2}\cdot 2\mu\,\Phi_{t}
= \big(1-\mu\eta\big)\Phi_{t}. \tag{14}
\]

\paragraph{[Step 12] Unroll the recursion over $m$ inner steps.}
Applying (14) repeatedly,
\[
\Phi_{m}\le (1-\mu\eta)^{m}\Phi_{0}.
\]
Recall $\Phi_{0}=f(\tilde{x}^{k})-f^{\star}$ because $x_{0}=\tilde{x}^{k}$ and $\|x_{0}-\tilde{x}^{k}\|=0$.
Hence
\[
\mathbb{E}\!\big[f(\tilde{x}^{k+1})-f^{\star}\big]=\mathbb{E}\!\big[f(x_{m})-f^{\star}\big]\le (1-\mu\eta)^{m}\big(f(\tilde{x}^{k})-f^{\star}\big). \tag{15}
\]

\paragraph{[Step 13] Choose $m$ to obtain a simple contraction per epoch.}
With $m\ge 2L/\mu$ and $\eta\le 1/(3L)$ we have
\[
(1-\mu\eta)^{m}\le \exp(-\mu\eta m)\le \exp\!\big(-\tfrac{2L}{\mu}\cdot\mu\eta\big)
= \exp(-2L\eta)\le 1-\tfrac{\mu\eta}{2},
\]
where the last inequality follows from the elementary bound $e^{-a}\le 1-a/2$ for $a\in[0,1]$,
and $L\eta\le 1/3$ guarantees $a=2L\eta\le 2/3<1$.
Thus (15) yields
\[
\mathbb{E}\!\big[f(\tilde{x}^{k+1})-f^{\star}\big]
\le\bigl(1-\tfrac{\mu\eta}{2}\bigr)\bigl[f(\tilde{x}^{k})-f^{\star}\bigr].
\]
Iterating over $k$ proves the claimed linear rate. \qedhere
\end{proof}

\section*{Rate Derivation (With Constants)}
From Theorem \ref{thm:linear} and the choice $\eta=1/(3L)$,
the contraction factor becomes $1-\mu/(6L)$. After $K$ epochs,
\[
\mathbb{E}\!\big[f(\tilde{x}^{K})-f^{\star}\big]
\le \Bigl(1-\frac{\mu}{6L}\Bigr)^{K}\bigl[f(\tilde{x}^{0})-f^{\star}\bigr].
\]
Each epoch costs $n$ gradient evaluations for the full gradient plus $m$ stochastic
gradients. With $m=2L/\mu$ the total number of stochastic gradient calls needed to achieve
$\epsilon$ accuracy is
\[
\mathcal{O}\!\Bigl(\bigl(n+\tfrac{2L}{\mu}\bigr)\log\frac{1}{\epsilon}\Bigr).
\]

\section*{Special Cases}
\begin{enumerate}
\item \textbf{Strongly convex $f$.} Strong convexity with parameter $\lambda$ implies the PL condition with $\mu=\lambda$; thus the theorem recovers the classical linear convergence of SVRG for strongly convex objectives.
\item \textbf{Non‑convex PL functions.} Even when $f$ is non‑convex but satisfies the PL inequality (e.g., certain over‑parameterised neural networks), the same linear rate holds, showing the broad applicability of SVRG‑PL.
\item \textbf{Limit $\mu\to0$.} If the PL constant tends to zero, the contraction factor approaches $1$, and the bound degrades to the sub‑linear rates typical for general smooth non‑convex optimization, consistent with known lower bounds.
\end{enumerate}

\end{document}